\section{Evaluation Methodology}\label{appendix:methodology}

We propose the following methodology for evaluating MARL methods using SMAC. 

To ensure the fairness of the challenge and comparability of results, performances should be evaluated under standardised conditions. 
One should not undertake any changes to the environment used for evaluating the policies. 
This includes the observation and state spaces, action space, the game mechanics, and settings of the environment (e.g., frame-skipping rate).
One should not modify the StarCraft II map files in any way or change the difficulty of the game AI. Episode limits of each scenario should also remain unchanged. %

SMAC restricts the execution of the trained models to be decentralised, i.e., during testing each agent must base its policy solely on its own action-observation history and cannot use the global state or the observations of other agents. %
It is, however, acceptable to train the decentralised policies in centralised fashion. Specifically, agents can exchange individual observations, model parameters and gradients during training as well as make use of the global state. %

\subsection{Evaluation Metrics}


Our main evaluation metric is the mean win percentage of evaluation episodes as a function of environment steps observed, over the course of training. 
Such progress can be estimated by periodically running a fixed number of evaluation episodes (in practice, 32) with any exploratory behaviours disabled. 
Each experiment is repeated using a number of independent training runs and the resulting plots include the median performance as well as the 25-75\% percentiles. 
We use five independent runs for this purpose in order to strike a balance between statistical significance and the computational requirements.
We recommend using the median instead of the mean in order to avoid the effect of any outliers. 
We report the number of independent runs, as well as environment steps used in training. %
Each independent run takes between 8 to 16 hours, depending on the exact scenario, using Nvidia Geforce GTX 1080 Ti graphics cards.



It can prove helpful to other researchers to include the computational resources used, and the wall clock time for running each experiment. SMAC provides functionality for saving StarCraft II replays, which can be viewed using a freely available client. The resulting videos can be used to comment on interesting behaviours observed.
